% 13-special-paths.tex — 特殊路径：Bayer、Tetrix 与 Profile 约束
% Source: chapters/13-special-paths-bayer-tetrix-and-profile-constraints.md

\chapter{特殊路径：Bayer、Tetrix 与 Profile 约束}\label{ch:special-paths}

\section{本章目标}

\begin{itemize}
  \item 理解 Bayer CFA 图像在 JPEG XS 中的处理方式
  \item 掌握 Star-Tetrix 变换的 4 遍扫描机制
  \item 理解 Profile/Level/Sublevel 三层约束体系
  \item 了解 MLS（Mathematically Lossless）和 DWT 抑制路径
\end{itemize}

\section{背景与标准语义}

JPEG XS 支持多种特殊输入路径，其中最重要的是 \textbf{Bayer CFA}（Color Filter Array）图像。Bayer 图像每个像素只有一个颜色分量（R/G/B 之一），通过 CFA 模式（RGGB/BGGR/GRBG/GBRG）排列。JPEG XS 使用 \textbf{Star-Tetrix} 变换将 Bayer 数据分解为 4 个虚拟分量，再分别编码。

Profile/Level/Sublevel 三层约束体系确保不同实现之间的互操作性。

\section{Bayer 到 4 分量的预处理}

\subsection{\texttt{mct\_bayer\_to\_4()}}

Bayer 输入是单分量图像，尺寸为 $W \times H$。预处理将其转换为 4 分量图像，每个分量尺寸为 $W/2 \times H/2$。

\begin{lstlisting}[numbers=none]
RGGB 模式下的 2x2 块：
+----+----+
| R  | G1 |   ->  comp[0]=R,  comp[1]=G1
+----+----+   ->  comp[2]=B,  comp[3]=G2
| G2 | B  |
+----+----+
\end{lstlisting}

代码中（\texttt{mct.c:413}）：

\begin{lstlisting}[language=C,numbers=none]
bool mct_bayer_to_4(xs_image_t* im, const xs_cfa_pattern_t cfa_pattern)
{
    // 输入: ncomps=1, width=W, height=H
    // 输出: ncomps=4, width=W/2, height=H/2

    for (int line = 0; line < im->height; ++line)
        for (int x = 0; x < im->width; ++x)
        {
            dst[0] = src[0];       // R  (左上)
            dst[2] = src[w];       // B  (左下)
            dst[1] = src[1];       // G1 (右上)
            dst[3] = src[w + 1];   // G2 (右下)
            src += 2;
        }
}
\end{lstlisting}

CFA 模式决定了 4 个虚拟分量到原始 Bayer 位置的映射：

\begin{table}[H]
\centering
\caption{CFA 模式与分量映射}
\begin{tabular}{@{}ccccc@{}}
\toprule
CFA 模式 & comp[0] & comp[1] & comp[2] & comp[3] \\
\midrule
RGGB & R & G1 & B & G2 \\
BGGR & B & G1 & R & G2 \\
GRBG & G1 & R & G2 & B \\
GBRG & G1 & B & G2 & R \\
\bottomrule
\end{tabular}
\end{table}

\subsection{逆过程 \texttt{mct\_4\_to\_bayer()}}

解码端将 4 分量合并回单分量 Bayer 图像（\texttt{mct.c:468}）。

\section{Star-Tetrix 变换}

\subsection{变换概述}

Star-Tetrix 是专为 Bayer CFA 设计的可逆色彩变换。它包含 \textbf{4 个遍历}（pass），每个遍历处理不同方向的邻域关系：

\begin{engineeringnote}
以下 4 遍扫描的功能描述（``利用水平/垂直相邻的同色像素''等）是对 \texttt{mct\_forward\_tetrix()} 中 4 个循环体的功能性解释。标准文本（ISO/IEC 21122）定义了每遍的精确数学公式，但未解释为何选择这 4 种特定的邻域方向。``去相关''是基于信号处理原理的推断。
\end{engineeringnote}

\begin{enumerate}
  \item \textbf{Pass 1 --- Forward CbCr}：利用水平/垂直相邻的同色像素计算色度残差
  \item \textbf{Pass 2 --- Forward Y}：利用色度信息预测亮度
  \item \textbf{Pass 3 --- Forward delta}：利用对角相邻像素计算残差
  \item \textbf{Pass 4 --- Forward average}：利用对角平均值进一步去相关
\end{enumerate}

\subsection{\texttt{mct\_tetrix\_access()} --- 像素寻址}

Tetrix 的核心难点是\textbf{跨分量像素寻址}：给定当前坐标 $(x, y)$、分量 $c$、偏移 $(r_x, r_y)$，需要找到对应的像素位置。

代码中（\texttt{mct.c:128}）：

\begin{lstlisting}[language=C,numbers=none]
xs_data_in_t* mct_tetrix_access(xs_image_t* im, const int c, const int Cf,
                                 const int Ct, const int rx, const int ry,
                                 const int x, const int y)
{
    int t_x = rx + ((Ct + c) & 1);
    int t_y = ry + (((~(c)) >> 1) & 1);
    const int k = ((((~(t_y)) << 1) & 2) | (((Ct) ^ (t_x)) & 1));
    // ... 边界处理和坐标变换
    return im->comps_array[k] + t_y * im->width + t_x;
}
\end{lstlisting}

参数说明：
\begin{itemize}
  \item $C_f$：Full（3）或 Inline（0）模式
  \item $C_t$：由 CFA 模式决定（RGGB/BGGR=0，GRBG/GBRG=1）
  \item $c$：目标分量（0--3）
\end{itemize}

\subsection{Forward CbCr（Pass 1）}

对分量 1 和 2（色度），利用水平相邻的同色像素：
\begin{equation}
s_c = s_c - \left\lfloor \frac{s_{c,\mathit{left}} + s_{c,\mathit{right}} + s_{c,\mathit{top}} + s_{c,\mathit{bottom}}}{4} \right\rfloor, \quad c \in \{1, 2\}
\end{equation}

\subsection{Forward Y（Pass 2）}

对分量 0 和 3（亮度），利用色度信息：
\begin{align}
s_0 &= s_0 + \left\lfloor \frac{(s_{0,\mathit{left}} + s_{0,\mathit{right}}) \ll e_2 + (s_{0,\mathit{top}} + s_{0,\mathit{bottom}}) \ll e_1}{8} \right\rfloor \\
s_3 &= s_3 + \left\lfloor \frac{(s_{3,\mathit{top}} + s_{3,\mathit{bottom}}) \ll e_2 + (s_{3,\mathit{left}} + s_{3,\mathit{right}}) \ll e_1}{8} \right\rfloor
\end{align}

其中 $e_1, e_2$ 是 CTS marker 中的参数（0--3），控制水平和垂直方向的加权。

\subsection{Forward delta（Pass 3）}

\begin{equation}
s_3 = s_3 - \left\lfloor \frac{s_{3,TL} + s_{3,TR} + s_{3,BL} + s_{3,BR}}{4} \right\rfloor
\end{equation}

TL/TR/BL/BR 是对角相邻的同分量像素。

\subsection{Forward average（Pass 4）}

\begin{equation}
s_0 = s_0 + \left\lfloor \frac{s_{0,TL} + s_{0,TR} + s_{0,BL} + s_{0,BR}}{8} \right\rfloor
\end{equation}

逆变换按完全相反的顺序执行：Pass 4 $\to$ Pass 3 $\to$ Pass 2 $\to$ Pass 1，每步取反操作。

\subsection{CTS 参数}

CTS marker（0xFF18）定义 Tetrix 参数：

\begin{table}[H]
\centering
\caption{CTS 参数}
\begin{tabular}{@{}cll@{}}
\toprule
参数 & 含义 & 取值 \\
\midrule
$C_f$ & Full/Inline 模式 & 0 = Inline, 3 = Full \\
$e_1$ & 水平方向加权移位 & 0--3 \\
$e_2$ & 垂直方向加权移位 & 0--3 \\
\bottomrule
\end{tabular}
\end{table}

\section{DWT 抑制（$S_d$）}

当 $S_d > 0$ 时，最后 $S_d$ 个分量\textbf{不做 DWT}，直接以原始分辨率编码。典型用于 Bayer 路径：$S_d = 1$，最后一个分量（高分辨率亮度）跳过 DWT。

验证规则（\texttt{xs\_config.c:836}）：
\begin{lstlisting}[language=C,numbers=none]
if (cfg->p.Sd != 1)
{
    fprintf(stderr, "Error: Sd must be 1 for Bayer profiles\n");
    return false;
}
\end{lstlisting}

\section{Profile/Level/Sublevel 约束}

\subsection{Profile 约束表}

Profile 定义了编码工具集和参数范围。下表各列含义：

\begin{itemize}
  \item \textbf{颜色变换}：NONE=禁止，AUTO=自动选择（RCT 或 NONE），TETRIX=Bayer 专用变换
  \item \textbf{NLx}：最大水平 DWT 分解层数（0 表示不做 DWT）
  \item \textbf{Nbpp}：Sublevel FULL 时的最大 bpp（bits per pixel），即该 Profile 允许的最高码率
  \item \textbf{特殊约束}：每个 Profile 特有的额外限制
\end{itemize}

\begin{table}[H]
\centering
\caption{Profile 约束表}
\small
\begin{tabular}{@{}L{2.6cm}L{1.4cm}L{0.8cm}L{0.8cm}L{4.5cm}@{}}
\toprule
Profile & 颜色变换 & NLx & Nbpp & 特殊约束 \\
\midrule
LIGHT\_422\_10 & NONE & 1 & 20 & 4:2:2, 10-bit \\
LIGHT\_444\_12 & AUTO & 1 & 36 & 4:4:4, 12-bit \\
MAIN\_420\_12 & NONE & 1 & 18 & 4:2:0, 12-bit \\
MAIN\_444\_12 & AUTO & 1 & 36 & 4:4:4, 12-bit \\
HIGH\_420\_12 & NONE & 2 & 18 & 4:2:0, 12-bit, $NL_x=2$（更多分解层） \\
HIGH\_444\_12 & AUTO & 2 & 36 & 4:4:4, 12-bit, $NL_x=2$ \\
LIGHT\_BAYER & TETRIX & 0 & 64 & $S_d=1$（$NL_x=0$，不做 DWT 分解） \\
MAIN\_BAYER & TETRIX & 1 & 64 & $S_d=1$（最后一个分量跳过 DWT） \\
HIGH\_BAYER & TETRIX & 2 & 64 & $S_d=1$, $NL_x=2$（更多 DWT 分解层） \\
MLS\_12 & AUTO & 2 & 48 & $F_q=0$, $B_w=d$（数学无损，逐 bit 精确） \\
UNRESTRICTED & AUTO & --- & --- & 无任何约束，仅开发调试用 \\
\bottomrule
\end{tabular}
\end{table}

\subsection{Level 约束}

Level 定义图像尺寸上限：

\begin{table}[H]
\centering
\caption{Level 约束表}
\small
\begin{tabular}{@{}lrrrr@{}}
\toprule
Level & $W_{\max}$ & $H_{\max}$ & $L_{\max}$（像素） & $R_{s,\max}$（采样率） \\
\midrule
1K\_1 & 1280 & 5120 & 2,621,440 & 83,558,400 \\
2K\_1 & 2048 & 8192 & 4,194,304 & 133,693,440 \\
4K\_1 & 4096 & 16384 & 8,912,896 & 267,386,880 \\
8K\_1 & 8192 & 32768 & 35,651,584 & 1,069,547,520 \\
10K\_1 & 10240 & 40960 & 104,857,600 & 3,342,336,000 \\
\bottomrule
\end{tabular}
\end{table}

\subsection{Sublevel 约束}

Sublevel 定义最大码率：

\begin{table}[H]
\centering
\caption{Sublevel 约束表}
\begin{tabular}{@{}lc@{}}
\toprule
Sublevel & 最大 bpp \\
\midrule
2\_BPP & 2 \\
3\_BPP & 3 \\
6\_BPP & 6 \\
9\_BPP & 9 \\
12\_BPP & 12 \\
FULL & Profile 定义的 Nbpp \\
UNRESTRICTED & 无限制 \\
\bottomrule
\end{tabular}
\end{table}

\subsection{自动选择流程}

\texttt{xs\_config\_resolve\_auto\_values()}（\texttt{xs\_config.c:487}）：

\begin{lstlisting}[numbers=none]
1. Profile AUTO -> 根据 ncomps 和 subsampling 选择
2. Level AUTO -> 选择满足 WxH 的最低 Level
3. Sublevel AUTO -> 根据 target_bpp 选择最低 Sublevel
4. Bw/Fq -> 根据 Profile 和 NLT 解析
5. CAP bits -> 根据使用的功能自动设置
6. Gains/Priorities -> 查 LUT 或计算
\end{lstlisting}

\subsection{验证逻辑}

\texttt{xs\_config\_validate()} 执行以下检查：
\begin{itemize}
  \item 图像尺寸与 Level 匹配
  \item 码率与 Sublevel 匹配
  \item 颜色变换与分量格式匹配（RCT 需要 3 分量无下采样）
  \item NLy 约束：$\max(s_y) - 1 \leq NL_y \leq NL_x$
  \item $S_d$ 约束：$S_d < N_c$，被抑制的分量不能下采样
  \item Bayer Profile 必须使用 Tetrix
\end{itemize}

\section{MLS（数学无损）路径}

MLS Profile（\texttt{XS\_PROFILE\_MLS\_12}）的特点：
\begin{itemize}
  \item $F_q = 0$：无小数位，不做右移
  \item $B_w = \mathit{depth}$：工作位宽等于输入位深
  \item 禁用 rate control：\texttt{bitstream\_size\_in\_bytes = (size\_t)-1}
  \item 无 NLT
\end{itemize}

这确保所有变换都是整数可逆的，编码结果数学上完全等于输入。

\subsection{流程图}

\begin{figure}[H]
\centering
\begin{tikzpicture}[
    node distance=0.8cm,
    block/.style={rectangle, draw, fill=blue!10, text width=3.2cm, text centered, minimum height=0.8cm, font=\small},
    decision/.style={diamond, draw, fill=yellow!10, text width=2cm, text centered, aspect=2, font=\small},
    arrow/.style={-{Stealth[length=2.5mm]}, thick}
]
    \node[block] (bayer) {Bayer CFA 图像};
    \node[block, below=of bayer] (b24) {\texttt{mct\_bayer\_to\_4}\\单分量$\to$4 分量};
    \node[block, below=of b24] (tetrix) {\texttt{mct\_forward\_tetrix}\\4 遍扫描};
    \node[block, below=of tetrix] (dwt) {DWT 正变换};
    \node[block, below=of dwt] (enc) {常规编码管线};
    \node[block, below=of enc] (stream) {JPEG XS 码流};

    \node[block, right=3cm of bayer] (rgb) {常规 RGB/YUV};
    \node[block, below=of rgb] (nlt) {NLT 正变换};
    \node[block, below=of nlt] (mct) {MCT: none/RCT};

    \draw[arrow] (bayer) -- (b24);
    \draw[arrow] (b24) -- (tetrix);
    \draw[arrow] (tetrix) -- (dwt);
    \draw[arrow] (dwt) -- (enc);
    \draw[arrow] (enc) -- (stream);

    \draw[arrow] (rgb) -- (nlt);
    \draw[arrow] (nlt) -- (mct);
    \draw[arrow] (mct.south) -- ++(0,-0.6) -| (dwt.east);
\end{tikzpicture}
\caption{Bayer 路径与常规路径的流程图}
\label{fig:special-paths-flowchart}
\end{figure}

\section{实现映射}

\begin{implmap}
\begin{tabular}{@{}L{2.8cm}L{3.3cm}L{5.4cm}L{4.0cm}@{}}
\toprule
标准概念 & 入口文件 & 关键函数/结构 & 说明 \\
\midrule
Bayer$\to$4 & \btt{mct.c:413} & \btt{mct\_bayer\_to\_4()} & 单分量$\to$4 分量 \\
4$\to$Bayer & \btt{mct.c:468} & \btt{mct\_4\_to\_bayer()} & 4 分量$\to$单分量 \\
Tetrix 正变换 & \btt{mct.c:157} & \btt{mct\_forward\_tetrix()} & 4 遍扫描 \\
Tetrix 逆变换 & \btt{mct.c:246} & \btt{mct\_inverse\_tetrix()} & 逆序 4 遍 \\
像素寻址 & \btt{mct.c:128} & \btt{mct\_tetrix\_access()} & 跨分量寻址 \\
Profile 表 & \btt{xs\_config.c:91} & \btt{XS\_PROFILES[]} & 15 个 Profile \\
Level 表 & \btt{xs\_config.c:113} & \btt{XS\_LEVELS[]} & 10 个 Level \\
Sublevel 表 & \btt{xs\_config.c:127} & \btt{XS\_SUBLEVELS[]} & 7 个 Sublevel \\
配置解析 & \btt{xs\_config.c:487} & \btt{xs\_config\_resolve\_auto\_values()} & AUTO 解析 \\
配置验证 & \btt{xs\_config.c:567} & \btt{xs\_config\_validate()} & 约束检查 \\
\bottomrule
\end{tabular}
\end{implmap}

\section{常见误解}

\begin{misunderstanding}
\begin{enumerate}
  \item \textbf{``Bayer 图像直接 DWT''} --- 不是。Bayer 先通过 \texttt{mct\_bayer\_to\_4()} 转换为 4 分量，再做 MCT 和 DWT。
  \item \textbf{``Star-Tetrix 和 RCT 一样是简单的线性变换''} --- 不一样。Tetrix 是 4 遍扫描，每遍使用不同方向的邻域，且包含跨分量寻址。
  \item \textbf{``Profile 只影响编码参数''} --- Profile 还约束了可用的编码工具。例如非 Bayer Profile 不允许 NLT，非 UNRESTRICTED Profile 要求 \texttt{slice\_height=16}。
  \item \textbf{``Level 只约束图像尺寸''} --- Level 还约束了最大采样率 $R_{s,\max}$，限制了帧率。
\end{enumerate}
\end{misunderstanding}

\section{本章小结}

JPEG XS 通过 Bayer$\to$4 分量$\to$Tetrix 路径支持 CFA 图像。Star-Tetrix 是 4 遍扫描的可逆色彩变换，利用 Bayer 模式的空间规律去相关。Profile/Level/Sublevel 三层约束体系确保互操作性，自动解析流程简化了用户配置。

下一章将通过一个端到端的编码追踪示例，将所有步骤串联起来。
